\subsection{Purpose}

The purpose of this section is to establish minimum requirements for designing, configuring, operating, securing, monitoring, and maintaining the organization's network infrastructure.

These requirements are intended to:

\begin{itemize}
    \item Protect network services and systems from unauthorized access
    \item Support the confidentiality, integrity, and availability of systems and data
    \item Ensure that network devices and connections are securely configured and maintained
    \item Limit network access to authorized users, devices, systems, and services
    \item Reduce the likelihood and impact of network-based security incidents
    \item Establish requirements for network segmentation, remote access, wireless access, and firewall management
    \item Ensure that network changes are authorized, documented, tested, and recoverable
    \item Support reliable operation, troubleshooting, incident response, and service restoration
\end{itemize}

Network infrastructure \must{} be managed to protect the confidentiality, integrity, and availability of systems and data. Network devices, services, and connections \must{} have a documented business purpose, designated owner, and appropriate security configuration.

\subsection{Network Inventory and Ownership}

All network infrastructure \must{} be recorded in the asset inventory. This includes firewalls, routers, switches, wireless access points, VPN systems, network appliances, virtual networks, cloud networks, and externally accessible services.

The inventory \should{} include:
\begin{itemize}
    \item Device or service name
    \item Device type and physical or logical location
    \item Manufacturer, model, and serial number, where applicable
    \item IP addresses, hostnames, and network segments
    \item System or service owner
    \item Technical administrator
    \item Business purpose
    \item Internet exposure and external dependencies
    \item Firmware or operating system version
    \item Support status and expected replacement date
    \item Date of the last review
\end{itemize}

Each significant network segment, connection, firewall, and externally accessible service \must{} have a designated owner.

\subsection{Network Architecture and Segmentation}

Networks \must{} be designed and configured according to business and security requirements. Systems with different levels of trust or sensitivity \should{} be separated using appropriate controls such as VLANs, routing rules, firewalls, security groups, or separate physical networks.

Where technically and operationally appropriate:
\begin{itemize}
    \item User workstations \should{} be separated from server and infrastructure networks
    \item Guest and personal devices \must{} be isolated from internal business systems
    \item Management interfaces \must{} be restricted to authorized administrators
    \item Internet-facing systems \must{} be separated from internal networks
    \item Backup networks and storage \should{} be protected from unauthorized access and routine user traffic
    \item Development, testing, and production environments \should{} be separated
    \item Systems handling sensitive information \should{} be placed in appropriately protected network segments
\end{itemize}

Network access \must{} be denied by default unless it is explicitly required and authorized. Firewall and access-control rules \must{} specify the source, destination, service, purpose, and owner.

\subsection{Firewalls and Access Controls}

Firewalls, routers, wireless systems, VPNs, and other network security devices \must{} use secure configurations.

Network access rules:
\begin{itemize}
    \item \must{} be limited to the minimum required protocols, ports, addresses, and users
    \item \must{} have a documented business or technical purpose
    \item \must{} be reviewed periodically
    \item \must{} be removed when no longer required
    \item \must{} be tested after significant changes
    \item \should{} include an expiration date for temporary access
\end{itemize}

Broad rules that allow unrestricted access between networks \should{} be avoided. Rules that permit access from the Internet to internal systems \must{} be specifically approved and protected with appropriate authentication and monitoring.

Administrative access to network devices \must{} use individual administrator accounts where supported. Shared administrative accounts \should{} be avoided and, when required, \must{} be documented and controlled.

\subsection{Wireless Networks}

Wireless networks \must{} use security settings appropriate to the capabilities of the devices and the sensitivity of the information being transmitted.

Wireless networks:
\begin{itemize}
    \item \must{} use strong encryption and authentication
    \item \must{} use separate networks for guests and internal users where guest access is provided
    \item \must{} restrict administrative access to authorized personnel
    \item \must{} use a unique administrative password or other appropriate administrative authentication
    \item \should{} use enterprise authentication where supported
    \item \mustnot{} use default administrator passwords or default security settings
\end{itemize}

Guest wireless access \must{} not provide access to internal systems unless the access has been specifically authorized and secured.

\subsection{Remote Access and VPN}

Remote access to internal systems \must{} be authorized, encrypted, and limited to the systems and services required for the approved purpose.

Remote access systems \must{}:
\begin{itemize}
    \item Require individual user accounts
    \item Use MFA where supported
    \item Use approved encryption protocols
    \item Enforce access restrictions based on user, device, network, or other appropriate conditions
    \item Log successful and unsuccessful authentication attempts
    \item Remove or disable access when it is no longer required
\end{itemize}

Remote administrative access \must{} be restricted to authorized administrators and \should{} be limited to approved devices and managed network connections. Remote access tools that bypass organizational security controls \mustnot{} be installed or used without approval.

Third-party remote access \must{} be limited to the required systems, time period, and permissions. Temporary access \should{} expire automatically whenever technically possible.

\subsection{Network Device Configuration}

Network devices \must{} be securely configured before being placed into production.

Administrators \must{}:
\begin{itemize}
    \item Change default passwords and disable unused default accounts
    \item Disable unnecessary services, interfaces, and protocols
    \item Apply approved firmware and security updates
    \item Use secure management protocols
    \item Restrict management access to authorized networks or administrator devices
    \item Synchronize device time with an approved time source
    \item Save and protect current configurations
    \item Document significant configuration changes
    \item Back up configurations when needed for recovery
\end{itemize}

Insecure protocols such as unencrypted administrative services \mustnot{} be used where a secure alternative is available. Network configuration backups \must{} be protected from unauthorized access because they may contain credentials, keys, addresses, or other sensitive information.

\subsection{Network Changes}

Changes to network infrastructure, firewall rules, routing, wireless configuration, VPN settings, or network security controls \must{} follow the approved change-management process.

A significant network change \must{} include:
\begin{itemize}
    \item A description and business purpose
    \item The affected systems and services
    \item Security and availability considerations
    \item Implementation steps
    \item Testing or validation steps
    \item A rollback or recovery plan
    \item The responsible administrator
    \item Required approval
    \item The date of implementation
\end{itemize}

Emergency network changes may be implemented when necessary to protect systems or restore service, but they \must{} be documented and reviewed afterward.

\subsection{Network Availability and Capacity}

Network infrastructure \must{} be monitored and maintained to support required business operations.

Administrators \should{} monitor:
\begin{itemize}
    \item Availability and service interruptions
    \item Bandwidth and utilization
    \item Capacity and performance
    \item Hardware and interface errors
    \item Unusual traffic patterns
    \item VPN and wireless authentication failures
    \item Firewall and intrusion-prevention alerts
    \item Expiring certificates and service credentials
\end{itemize}

Critical network components \should{} have appropriate redundancy, documented recovery procedures, and tested restoration procedures based on their business importance.

\subsection{Network Security Monitoring}

Network devices and services \must{} generate security and operational logs appropriate to their function. Relevant logs \should{} be forwarded to a centralized logging system where technically and operationally feasible.

Network logs \should{} include:
\begin{itemize}
    \item Administrative logins and configuration changes
    \item Firewall rule changes
    \item VPN connections and disconnections
    \item Authentication failures
    \item Blocked and permitted security events
    \item Significant routing or interface changes
    \item Wireless access and authentication events
    \item Malware, intrusion, or denial-of-service alerts
\end{itemize}

Network logs \mustnot{} contain passwords, private keys, recovery secrets, or other sensitive credential values.

\subsection{Network Review}

Network infrastructure and access rules \must{} be reviewed periodically and after significant organizational, technical, or security changes.

The review \should{} verify that:
\begin{itemize}
    \item Network devices and connections are accurately inventoried
    \item Owners and administrators are assigned
    \item Firewall and access-control rules remain necessary
    \item Unused services, ports, accounts, and connections have been removed
    \item Remote access remains authorized
    \item Wireless networks remain appropriately secured
    \item Firmware and software remain supported
    \item Network diagrams and documentation are current
    \item Security and availability issues have been assigned for remediation
\end{itemize}

Any network component that cannot be securely maintained \must{} be isolated, replaced, or formally approved as an exception.